
\documentclass{article}
\usepackage{colortbl}
\usepackage{makecell}
\usepackage{multirow}
\usepackage{supertabular}

\begin{document}

\newcounter{utterance}

\centering \large Interaction Transcript for game `cladder', experiment `full\_v1.5\_default', episode 2621 with selfrefine{-}Qwen3.5{-}9B.
\vspace{24pt}

{ \footnotesize  \setcounter{utterance}{1}
\setlength{\tabcolsep}{0pt}
\begin{supertabular}{c@{$\;$}|p{.15\linewidth}@{}p{.15\linewidth}p{.15\linewidth}p{.15\linewidth}p{.15\linewidth}p{.15\linewidth}}
    \# & \multicolumn{2}{c}{Player} && \multicolumn{2}{c}{Game Master} \\
    \hline

    \theutterance \stepcounter{utterance}  
    & & & \multicolumn{4}{p{0.6\linewidth}}{
        \cellcolor[rgb]{0.9,0.9,0.9}{
            \makecell[{{p{\linewidth}}}]{
                \texttt{\tiny{[P1$\langle$GM]}}
                \texttt{You are an expert in causal inference. The following question is not a typical commonsense query, but rather a meticulously designed question created by a professor specializing in causal inference, intended to assess the students' mastery of the course content.} \\
\texttt{Imagine a self{-}contained, hypothetical world with only the following conditions, and without any unmentioned factors or causal relationships: Gender has a direct effect on smoking and lung cancer. Smoking has a direct effect on tar deposit. Tar deposit has a direct effect on lung cancer. Gender is unobserved. Method 1: We look directly at how smoking correlates with lung cancer in general. Method 2: We look at this correlation case by case according to gender. To understand how smoking affects lung cancer, is it more correct to use the Method 1 than Method 2?} \\
\texttt{Start your answer with 'Yes' or 'No', followed by additional reasoning or evidence to support your explanation.} \\
            }
        }
    }
     \\ \\

    \theutterance \stepcounter{utterance}  
    & \multicolumn{4}{p{0.6\linewidth}}{
        \cellcolor[rgb]{0.561,0.737,0.561}{
            \makecell[{{p{\linewidth}}}]{
                \texttt{\tiny{[P1$\rangle$GM]}}
                \texttt{No} \\
\\ 
\texttt{**Reasoning:**} \\
\\ 
\texttt{In this hypothetical world, the causal structure forms a Directed Acyclic Graph (DAG) where **Gender** is a common cause (confounder) of both **Smoking** and **Lung Cancer**. Specifically:} \\
\texttt{1.  Gender \$\textbackslash{}rightarrow\$ Smoking} \\
\texttt{2.  Gender \$\textbackslash{}rightarrow\$ Lung Cancer} \\
\texttt{3.  Smoking \$\textbackslash{}rightarrow\$ Tar Deposit \$\textbackslash{}rightarrow\$ Lung Cancer} \\
\\ 
\texttt{To estimate the causal effect of smoking on lung cancer, we must isolate the pathway \$Smoking \textbackslash{}rightarrow Tar \textbackslash{}rightarrow Lung Cancer\$ from the non{-}causal (backdoor) pathway \$Smoking \textbackslash{}leftarrow Gender \textbackslash{}rightarrow Lung Cancer\$.} \\
\\ 
\texttt{*   **Method 1 (Crude Association):** This method calculates the marginal correlation between smoking and lung cancer (\$P(Lung Cancer | Smoking)\$). Because Gender is a confounder that influences both the exposure (smoking) and the outcome (lung cancer), it creates a spurious association. Even if smoking had no effect on lung cancer, the fact that certain genders smoke more and have higher baseline lung cancer rates would make smoking appear correlated with lung cancer. Therefore, Method 1 yields a biased estimate of the causal effect.} \\
\\ 
\texttt{*   **Method 2 (Stratification/Adjustment):** This method calculates the association within specific strata of Gender (e.g., comparing smokers and non{-}smokers separately for males and females). By conditioning on Gender, we block the backdoor path (\$Smoking \textbackslash{}leftarrow Gender \textbackslash{}rightarrow Lung Cancer\$). In each stratum, the distribution of Gender is fixed, eliminating its confounding influence. The resulting correlation within these strata reflects the true causal mechanism mediated by tar deposits.} \\
\\ 
\texttt{Since the goal is to understand "how smoking affects lung cancer" (the causal effect), Method 2 is the correct approach because it controls for the confounding variable (Gender), whereas Method 1 fails to do so.} \\
            }
        }
    }
    & & \\ \\

    \theutterance \stepcounter{utterance}  
    & & & \multicolumn{2}{p{0.3\linewidth}}{
        \cellcolor[rgb]{0.9,0.9,0.9}{
            \makecell[{{p{\linewidth}}}]{
                \texttt{\tiny{[GM$|$GM]}}
                \texttt{No} \\
            }
        }
    }
    & & \\ \\

    \theutterance \stepcounter{utterance}  
    & & & \multicolumn{2}{p{0.3\linewidth}}{
        \cellcolor[rgb]{0.9,0.9,0.9}{
            \makecell[{{p{\linewidth}}}]{
                \texttt{\tiny{[GM$|$GM]}}
                \texttt{no} \\
            }
        }
    }
    & & \\ \\

    \theutterance \stepcounter{utterance}  
    & & & \multicolumn{2}{p{0.3\linewidth}}{
        \cellcolor[rgb]{0.9,0.9,0.9}{
            \makecell[{{p{\linewidth}}}]{
                \texttt{\tiny{[GM$|$GM]}}
                \texttt{game\_result = WIN} \\
            }
        }
    }
    & & \\ \\

\end{supertabular}
}

\end{document}
